% Motor: XeLaTeX. Sesión 24 - material docente UTP.
% Estructura docente de la Sesión 23.
\documentclass[aspectratio=169,10pt,c]{beamer}
% Reutiliza sin cambios la identidad visual UTP validada en la Sesion 08.
% Reutiliza sin cambios la identidad visual UTP validada en la Sesion 05.
% La segunda ruta permite compartir los assets comunes sin duplicar binarios.
% !TEX root = slides.tex
% Preambulo comun: estilo UTP, colores, fuentes, plantillas de beamer y
% comandos de apoyo didactico. Se carga desde slides.tex tras \documentclass.

% La opción "c" de beamer reparte el espacio libre 1:1.5 (arriba:abajo), por lo
% que el contenido queda por encima del centro. Igualamos ambos rellenos.
\makeatletter
\define@key{beamerframe}{c}[true]{%
  \beamer@frametopskip=0pt plus 1fill\relax%
  \beamer@framebottomskip=0pt plus 1fill\relax%
  \beamer@frametopskipautobreak=0pt plus .5\paperheight\relax%
  \beamer@framebottomskipautobreak=0pt plus .5\paperheight\relax%
}
\makeatother

\usepackage{graphicx}
\usepackage{tikz}
\usepackage{xcolor}
\usepackage{iftex}

\ifPDFTeX
  \usepackage[utf8]{inputenc}
  \usepackage[T1]{fontenc}
  \usepackage{lmodern}          % tipografia escalable para el respaldo con pdflatex
  \renewcommand{\familydefault}{\sfdefault}
\else
  \usepackage{fontspec}
  \IfFontExistsTF{Arial}
    {\setsansfont{Arial}[BoldFont={Arial Bold}]}
    {\setsansfont{Helvetica}[BoldFont={Helvetica Bold}]}
\fi

\usepackage[spanish,es-nodecimaldot]{babel}
\usepackage{microtype}
\usepackage{booktabs}
\usepackage{array}
\usetikzlibrary{arrows.meta,positioning,fit,calc,shapes.geometric}

\graphicspath{{assets/}}

% ---------- Identidad visual UTP ----------
\definecolor{utpRojo}{HTML}{C42B35}
\definecolor{utpNegro}{HTML}{111111}
\definecolor{utpGris}{HTML}{5F5F5F}
\definecolor{utpGrisMedio}{HTML}{9A9A9A}
\definecolor{utpGrisClaro}{HTML}{F1F1F1}

\usetheme{default}
\setbeamercolor{normal text}{fg=utpNegro,bg=white}
\setbeamercolor{structure}{fg=utpRojo}
\setbeamercolor{frametitle}{fg=utpNegro,bg=white}
\setbeamercolor{title}{fg=utpNegro}
\setbeamercolor{subtitle}{fg=utpNegro}
\setbeamercolor{author}{fg=utpNegro}
\setbeamercolor{date}{fg=utpNegro}
\setbeamercolor{block title}{fg=white,bg=utpNegro}
\setbeamercolor{block body}{fg=utpNegro,bg=utpGrisClaro}
\setbeamercolor{alertblock title}{fg=white,bg=utpRojo}
\setbeamercolor{alertblock body}{fg=utpNegro,bg=utpRojo!7}
\setbeamercolor{exampleblock title}{fg=white,bg=utpGris}
\setbeamercolor{exampleblock body}{fg=utpNegro,bg=utpGris!8}
\setbeamercolor{terminal}{fg=white,bg=utpNegro}

\setbeamerfont{title}{series=\bfseries,size=\fontsize{19.5}{22.5}\selectfont}
\setbeamerfont{subtitle}{size=\fontsize{11.5}{13.5}\selectfont}
\setbeamerfont{date}{series=\bfseries,size=\fontsize{11.5}{13.5}\selectfont}
\setbeamerfont{author}{size=\fontsize{10.5}{12.5}\selectfont}
\setbeamerfont{frametitle}{series=\bfseries,size=\large}
\setbeamerfont{block title}{series=\bfseries,size=\small}
\setbeamerfont{block body}{size=\normalsize}

\setbeamertemplate{navigation symbols}{}
\setbeamertemplate{headline}{}
\setbeamertemplate{blocks}[default]
\setbeamertemplate{itemize item}{\color{utpRojo}\small\raise.15ex\hbox{$\blacktriangleright$}}
\setbeamertemplate{itemize subitem}{\color{utpGris}\scriptsize$\bullet$}
\setbeamertemplate{enumerate item}{\color{utpRojo}\insertenumlabel.}
\setbeamertemplate{background canvas}{\color{white}\rule{\paperwidth}{\paperheight}}

% Título de diapositiva: negrita negra a la izquierda, logo UTP a la derecha
\setbeamertemplate{frametitle}{%
  \begin{tikzpicture}[remember picture,overlay]
    \node[anchor=north east,inner sep=0pt]
      at ([xshift=-7mm,yshift=-5.5mm]current page.north east)
      {\includegraphics[width=17mm]{utp-logo.png}};
  \end{tikzpicture}%
  \vspace{1.2mm}%
  \begin{minipage}[t]{.80\textwidth}\raggedright\insertframetitle\end{minipage}\par
  \vspace{1.1mm}{\color{utpRojo}\rule{\textwidth}{1.1pt}}
}

% Solo numeración discreta; la plantilla UTP no lleva pie de página
\setbeamertemplate{footline}{%
  \hfill{\color{utpGrisMedio}\scriptsize\insertframenumber}\hspace{6mm}\vspace{2.5mm}%
}

% Portada UTP
\setbeamertemplate{title page}{%
  \begin{tikzpicture}[remember picture,overlay]
    \node[anchor=center,inner sep=0pt] at (current page.center) {%
      \includegraphics[width=\paperwidth,height=\paperheight]{utp-cover-bg.jpeg}%
    };
    \node[anchor=north,align=center,text=utpNegro,inner sep=0pt,
      text width=145mm,font=\usebeamerfont{title}]
      at ([yshift=-25.2mm]current page.north) {\inserttitle};
    \node[anchor=north,align=center,text=utpNegro,inner sep=0pt,
      text width=142mm,font=\usebeamerfont{subtitle}]
      at ([yshift=-38.2mm]current page.north) {\insertsubtitle};
    \node[anchor=north,align=center,text=utpNegro,inner sep=0pt,
      font=\usebeamerfont{date}]
      at ([yshift=-50.4mm]current page.north) {\insertdate};
    \node[anchor=north,align=center,text=utpNegro,inner sep=0pt,
      font=\usebeamerfont{author}]
      at ([yshift=-58.4mm]current page.north) {\insertauthor};
  \end{tikzpicture}%
}

% ---------- Comandos de apoyo didáctico ----------
\newcommand{\concepto}[2]{%
  \begin{block}{#1}\large #2\end{block}}
\newcommand{\pregunta}[1]{%
  \begin{alertblock}{Decide y explica}#1\end{alertblock}}
\newcommand{\analogia}[1]{%
  \begin{exampleblock}{En palabras simples}\small #1\end{exampleblock}}
\newcommand{\enlarealidad}[1]{%
  \begin{block}{Cómo se ve en la realidad}\small #1\end{block}}
\newcommand{\ojo}[1]{%
  \begin{alertblock}{Ojo con esto}\small #1\end{alertblock}}
\newcommand{\apoyo}{{\normalfont\footnotesize\color{utpGris}~· material de apoyo}}
\newcommand{\marca}[1]{\tikz[baseline=(m.base)]{%
  \node[rounded corners=1.5pt,fill=utpGrisClaro,draw=utpGrisMedio,inner xsep=3pt,inner ysep=1.5pt,
        font=\footnotesize] (m) {#1};}}
\newenvironment{terminal}{%
  \begin{beamercolorbox}[wd=\linewidth,sep=5pt,rounded=true]{terminal}%
  \ttfamily\small\ignorespaces}{%
  \end{beamercolorbox}}

\tikzset{
  caja/.style={draw=utpNegro, very thick, rounded corners=1.5pt, fill=white,
    align=center, minimum height=0.72cm, inner sep=5pt, font=\small},
  cajaazul/.style={caja, draw=utpGris, fill=white},
  cajaverde/.style={caja, draw=utpNegro, fill=utpGrisClaro},
  cajanaranja/.style={caja, draw=utpRojo!75, fill=utpRojo!7},
  flecha/.style={-{Stealth[length=2.3mm]}, very thick, draw=utpNegro},
  flechavuelta/.style={{Stealth[length=2.3mm]}-{Stealth[length=2.3mm]}, very thick, draw=utpNegro}
}

\graphicspath{{../../Sesion05/slides/assets/}}



\graphicspath{{assets/}{../../Sesion05/slides/assets/}}
\usepackage{colortbl}
\usepackage{listings}
\lstset{basicstyle=\ttfamily\footnotesize,backgroundcolor=\color{utpGrisClaro},
  frame=none,aboveskip=4pt,belowskip=4pt,showstringspaces=false,
  columns=fullflexible,keepspaces=true,breaklines=true,
  keywordstyle=\color{utpRojo}\bfseries,
  morekeywords={SELECT,FROM,WHERE,GROUP,BY,ORDER,DESC,LIMIT,AS,COUNT,SUM,ROUND,import},
  xleftmargin=5pt,xrightmargin=5pt,framesep=5pt}

\title{Integración de Servicios Cloud GCP}
\subtitle{Big Data y BigQuery: preguntas a millones de filas con SQL}
\date{Semana 12 - Sesión 24}
\author{Rel Guzman}

\newcommand{\separador}[2]{%
\begin{frame}[plain]
\vfill
\begin{center}
{\fontsize{29}{34}\selectfont\bfseries #1}\par
\vspace{2mm}
{\small #2}\par
\vspace{4mm}
{\color{utpRojo}\rule{8.5cm}{1.2pt}}
\end{center}
\vfill
\end{frame}}

% Tabla pequeña de pedidos que usan las diapositivas de SQL.
\newcommand{\filaA}{\rowcolor{utpRojo!30}}
\newcommand{\filaB}{\rowcolor{utpRojo!12}}
\newcommand{\filaC}{\rowcolor{utpGrisClaro}}

\begin{document}

{\setbeamertemplate{footline}{}\maketitle}

% 2
\begin{frame}{Revisión: Cloud SQL por dentro}
\begin{columns}[c,onlytextwidth]
\column{.50\textwidth}
\centering
\begin{tikzpicture}[font=\small]
  \draw[utpRojo,thick,rounded corners=5pt,fill=utpRojo!8] (0,0) rectangle (6.4,4.6);
  \node[anchor=north west,align=left] at (0.15,4.5) {\textbf{Instancia} \texttt{wordpress-db}\\\footnotesize Motor: MySQL 8.0};
  \draw[utpRojo,thick,rounded corners=4pt,fill=white] (0.4,0.3) rectangle (6.0,3.3);
  \node[anchor=north west] at (0.55,3.2) {\textbf{Base de datos} \texttt{wordpress}};
  \foreach \t/\x in {wp\_posts/1.45,wp\_users/3.2,wp\_comments/4.95} {
    \draw[utpRojo,fill=utpRojo!25] (\x-0.75,1.85) rectangle (\x+0.75,2.15);
    \draw[utpRojo] (\x-0.75,1.0) grid[xstep=0.5,ystep=0.283] (\x+0.75,1.85);
    \node[font=\scriptsize] at (\x,0.65) {\texttt{\t}};
  }
\end{tikzpicture}
\column{.46\textwidth}
\small
\textbf{Instancia:} el servidor MySQL que Google crea y mantiene. Al crearla eliges el motor.\par\medskip
\textbf{Base de datos:} un espacio con nombre dentro de la instancia. Una instancia puede tener varias.\par\medskip
\textbf{Tabla:} donde quedan las filas. WordPress creó las suyas al instalarse.\par\vspace{3mm}
\begin{alertblock}{Pregunta rápida}
\small Si mañana necesitas otra base de datos, ¿tienes que crear otra instancia?
\end{alertblock}
\end{columns}
\end{frame}

% 3
\begin{frame}{Caso de estudio: Pedidos Express crece}
\begin{block}{Enunciado}
\small Pedidos Express ya tiene cuatro áreas: \textbf{ventas}, \textbf{inventario}, \textbf{facturación} y \textbf{recursos humanos}. Cada una necesita guardar sus datos en Cloud SQL. Diseña cómo quedaría.
\end{block}
\begin{center}
\begin{tikzpicture}[font=\footnotesize,
 dato/.style={draw=utpRojo,thick,rounded corners=3pt,fill=white,minimum width=3.1cm,minimum height=0.8cm,align=center}]
  \node[dato] at (0,0.6) {Pedido P-1001 de Ana};
  \node[dato] at (3.5,0.6) {Stock de audífonos: 35};
  \node[dato] at (7.0,0.6) {Factura F-101: S/ 250};
  \node[dato] at (10.5,0.6) {Sueldo de un repartidor};
  \node[dato] at (1.75,-0.5) {Producto: Mouse, S/ 59};
  \node[dato] at (5.25,-0.5) {Pago de la factura F-101};
  \node[dato] at (8.75,-0.5) {Asistencia de un empleado};
\end{tikzpicture}
\end{center}
\begin{alertblock}{Responde}
\small 1. ¿Cuántas instancias creas? \quad 2. ¿Cuántas bases de datos y con qué nombres? \quad 3. ¿En qué base va cada dato de arriba?
\end{alertblock}
\end{frame}

% 4
\begin{frame}{Caso de estudio: una solución}
\begin{center}
\begin{tikzpicture}[font=\small,
 nivel/.style={draw=utpRojo,thick,rounded corners=4pt,align=center,minimum height=0.7cm},
 bd/.style={nivel,fill=utpRojo!10,minimum width=2.9cm},
 tabla/.style={font=\scriptsize\ttfamily,align=center}]
  \node[nivel,fill=utpRojo!35,minimum width=3.2cm] (c) at (0,2.6) {\textbf{Cloud SQL}};
  \node[nivel,fill=utpRojo!22,minimum width=5.6cm] (i) at (0,1.7) {Instancia: \texttt{mysql-server} \quad Motor: MySQL};
  \draw[-latex,thick] (c) -- (i);
  \foreach \n/\x/\t [count=\k] in {ventas/-5.1/{clientes\\pedidos},inventario/-1.7/{productos\\stock},facturacion/1.7/{facturas\\pagos},recursos\_humanos/5.1/{empleados\\asistencia}} {
    \node[bd] (b\k) at (\x,0.4) {BD: \texttt{\n}};
    \draw[-latex,thick] (i.south) -- (b\k.north);
    \node[tabla,below=1mm of b\k] {\t};
  }
\end{tikzpicture}
\end{center}
\vspace{-1mm}
\begin{columns}[T,onlytextwidth]
\column{.48\textwidth}
\begin{block}{¿Por qué así?}
\small \textbf{Una instancia:} un solo servidor que pagar y mantener; alcanza para empezar.\\
\textbf{Una base por área:} cada área encuentra rápido sus tablas y los sueldos no se mezclan con los pedidos.
\end{block}
\column{.48\textwidth}
\begin{exampleblock}{En la consola}
\small Cloud SQL $\rightarrow$ \texttt{mysql-server} $\rightarrow$ \textbf{Bases de datos} $\rightarrow$ \textbf{Crear base de datos}. Se repite una vez por área.
\end{exampleblock}
\end{columns}
\end{frame}

% 5
\begin{frame}{Revisión: la pregunta de hoy}
\begin{center}
\begin{tikzpicture}[font=\small,
 caja/.style={draw=utpRojo,thick,rounded corners=4pt,align=center}]
  \node[caja,fill=utpRojo!25,minimum width=2.2cm,minimum height=1cm] (app) at (0,0) {\textbf{Aplicación}};
  \node[caja,fill=white,minimum width=2.4cm,minimum height=1cm] (sql) at (4.4,0) {\textbf{Cloud SQL}\\\footnotesize BD \texttt{ventas}};
  \draw[-latex,thick] (app) -- node[above,font=\scriptsize,align=center]{guarda\\P-1001} node[below,font=\scriptsize,align=center]{1 fila,\\al instante} (sql);
  \foreach \i in {0,...,7} \draw[utpRojo,fill=white] (6.4+\i*0.12,-0.9+\i*0.12) rectangle (8.0+\i*0.12,0.1+\i*0.12);
  \node[font=\footnotesize,align=center,fill=white,inner sep=1pt] at (7.62,-0.05) {\textbf{1.2 millones}\\de pedidos};
  \node[font=\footnotesize] at (7.6,-1.3) {3 años de historial};
  \node[caja,fill=utpRojo!10,text width=2.9cm,font=\footnotesize] (g) at (11.3,0) {El gerente pregunta:\\\textbf{¿Qué categoría vendió más en 3 años?}};
  \draw[-latex,thick] (g.west) -- (9.1,0);
\end{tikzpicture}
\end{center}
\begin{columns}[T,onlytextwidth]
\column{.48\textwidth}
\begin{block}{Registrar}
\small La app guarda un pedido a la vez. Cada operación toca \textbf{una fila}. Cloud SQL está hecho para eso.
\end{block}
\column{.48\textwidth}
\begin{block}{Analizar}
\small La pregunta del gerente debe recorrer \textbf{todas las filas} del historial. Si corre en la misma base, la app se pone lenta mientras tanto.
\end{block}
\end{columns}
\par\vspace{2mm}
\begin{alertblock}{La idea de hoy}
\small Para preguntas sobre millones de filas se usa una herramienta pensada para analizar: \textbf{BigQuery}.
\end{alertblock}
\end{frame}

% 6
\begin{frame}[c]
\begin{block}{Logro de la sesión}
Al finalizar, podrás explicar qué es \textbf{Big Data} y por qué \textbf{BigQuery} responde rápido sobre millones de filas, cargar una tabla en BigQuery y responder preguntas de negocio con consultas \textbf{SQL} que filtran, agrupan y suman.
\end{block}
\end{frame}

% 7
\begin{frame}{Saberes previos}
\begin{itemize}\setlength\itemsep{11pt}
  \item ¿Cuántas filas crees que caben en una hoja de Excel?
  \item Pedidos Express recibe unos 1\,100 pedidos al día. ¿Cuántos pedidos tendrá después de 3 años? ¿Podrías abrirlos en Excel?
  \item Tienes la lista de todos los pedidos con su categoría y su monto. ¿Cómo sacarías cuánto se vendió en cada categoría?
\end{itemize}
\end{frame}

% 8
\begin{frame}[c]{Utilidad de esta sesión}
\begin{center}
\begin{minipage}{.84\textwidth}
{\large Las empresas acumulan años de ventas, cobros y registros de uso. Si sabes hacerles preguntas con \textbf{SQL} (qué se vende más, en qué mes, en qué distrito, en qué se gasta), esos datos se convierten en decisiones. Con \textbf{BigQuery} lo haces sin instalar ni administrar servidores.}
\end{minipage}
\end{center}
\end{frame}

% 9
\separador{Parte 1: ¿qué es Big Data?}{Cuando los datos ya no caben en las herramientas de siempre}

% 10
\begin{frame}{¿Cuándo los datos son ``grandes''?}
\begin{center}
\begin{tikzpicture}[font=\footnotesize,x=1.2cm]
  % Escala: cada marca es 10 veces la anterior (de 100 a 1000 millones).
  \foreach \e/\t in {0/100,1/1 mil,2/10 mil,3/100 mil,4/1 M,5/10 M,6/100 M,7/1000 M} {
    \draw[utpGrisMedio!60] (\e,0.4) -- (\e,3.5);
    \node[font=\scriptsize,color=utpGris] at (\e,0.2) {\t};
  }
  \foreach \y/\l/\v/\n in {3.1/Contactos de tu celular/0.699/500,
                          2.35/Tienda de barrio: 1 año/2.261/18\,250,
                          1.6/Pedidos Express: 3 años/4.079/1\,200\,000,
                          0.85/GPS de repartidores: 3 años/6.635/432\,000\,000} {
    \node[anchor=east,align=right] at (-0.15,\y) {\l};
    \fill[utpRojo!70] (0,\y-0.24) rectangle (\v,\y+0.24);
    \node[anchor=west,font=\scriptsize\bfseries] at (\v,\y) {\n};
  }
  \draw[utpNegro,very thick,dashed] (4.0206,0.4) -- (4.0206,3.5);
  \node[anchor=south,font=\scriptsize\bfseries,fill=white,inner sep=1pt] at (4.0206,3.5) {Límite de Excel: 1\,048\,576 filas};
\end{tikzpicture}
\end{center}
\vspace{-1mm}
{\scriptsize\color{utpGris} Cada línea vertical vale 10 veces la anterior (M = millones). GPS: 1 posición cada 5 s durante 30 min de reparto = 360 por pedido.}
\begin{block}{Big Data}
\small Datos tan grandes, que llegan tan rápido o tan variados que una hoja de cálculo o una sola computadora ya no los procesa a tiempo.
\end{block}
\end{frame}

% 11
\begin{frame}{Las tres V de Big Data}
\begin{center}
\begin{tikzpicture}[font=\small,
 v/.style={draw=utpRojo,thick,rounded corners=5pt,fill=white,text width=3.7cm,minimum height=3.5cm,align=center,anchor=north}]
  \node[v,fill=utpRojo!10] at (0,0) {{\Large\bfseries Volumen}\\[2pt]\footnotesize ¿Cuántos datos hay?\\[8pt]\small 1.2 millones de pedidos\\432 millones de posiciones GPS};
  \node[v] at (4.6,0) {{\Large\bfseries Velocidad}\\[2pt]\footnotesize ¿Qué tan rápido llegan?\\[8pt]\small Cada repartidor envía su ubicación cada 5 segundos, todo el día};
  \node[v,fill=utpRojo!10] at (9.2,0) {{\Large\bfseries Variedad}\\[2pt]\footnotesize ¿En qué forma llegan?\\[8pt]\small Tablas de pedidos, fotos de productos, reseñas escritas y coordenadas GPS};
\end{tikzpicture}
\end{center}
\begin{alertblock}{Basta con una}
\small Si solo una de las tres V crece mucho, la hoja de cálculo o la base de siempre empiezan a quedarse cortas.
\end{alertblock}
\end{frame}

% 12
\begin{frame}{Registrar y analizar no son lo mismo}
\begin{columns}[T,onlytextwidth]
\column{.48\textwidth}
{\color{utpRojo}\bfseries Registrar:} {\small ``Guarda el pedido P-1001 de Ana''}\par\vspace{2mm}
\centering\scriptsize\setlength\tabcolsep{4pt}
\begin{tabular}{@{}lllll@{}}
\toprule
\textbf{pedido} & \textbf{fecha} & \textbf{cliente} & \textbf{categoria} & \textbf{monto}\\
\midrule
P-0998 & 2025-12-31 & Luis & Moda & 98\\
P-0999 & 2025-12-31 & Rosa & Supermercado & 24\\
P-1000 & 2025-12-31 & Carlos & Electrónica & 59\\
\rowcolor{utpRojo!35} P-1001 & 2025-12-31 & Ana & Hogar & 149\\
\bottomrule
\end{tabular}\par\vspace{2mm}
\raggedright\footnotesize La app agrega \textbf{una fila nueva}, con todas sus columnas, y termina.
\column{.48\textwidth}
{\color{utpRojo}\bfseries Analizar:} {\small ``¿Cuánto se vendió por categoría?''}\par\vspace{2mm}
\centering\scriptsize\setlength\tabcolsep{4pt}
\begin{tabular}{@{}lll>{\columncolor{utpRojo!35}}l>{\columncolor{utpRojo!35}}l@{}}
\toprule
\textbf{pedido} & \textbf{fecha} & \textbf{cliente} & \textbf{categoria} & \textbf{monto}\\
\midrule
P-0001 & 2023-01-01 & Marta & Moda & 98\\
P-0002 & 2023-01-01 & Jorge & Electrónica & 295\\
$\vdots$ & $\vdots$ & $\vdots$ & $\vdots$ & $\vdots$\\
P-1001 & 2025-12-31 & Ana & Hogar & 149\\
\bottomrule
\end{tabular}\par\vspace{2mm}
\raggedright\footnotesize Lee \textbf{todas las filas} de los 3 años, pero solo \textbf{2 columnas}: \texttt{categoria} y \texttt{monto}. Suma los montos de cada categoría.
\end{columns}
\par\vspace{3mm}
\renewcommand{\arraystretch}{1.1}
\centering\footnotesize
\begin{tabular}{@{}p{3.6cm}>{\raggedright\arraybackslash}p{4.4cm}>{\raggedright\arraybackslash}p{4.4cm}@{}}
\toprule
 & \textbf{Registrar} {\color{utpGris}(Cloud SQL)} & \textbf{Analizar} {\color{utpGris}(BigQuery)}\\
\midrule
\textbf{Filas que usa} & 1, la del pedido nuevo & Todas: 1.2 millones\\
\textbf{Columnas que usa} & Todas las de esa fila & Solo las que pide la pregunta\\
\textbf{Cuánto tarda} & Milisegundos & Segundos\\
\textbf{Cuántas veces al día} & Miles: una por pedido & Pocas: una por reporte\\
\bottomrule
\end{tabular}
\end{frame}

% 13
\begin{frame}{¿Cómo responde rápido BigQuery? Dos ideas}
\begin{columns}[T,onlytextwidth]
\column{.47\textwidth}
\begin{block}{1. Guarda cada columna por separado}
\small Si la consulta usa \texttt{categoria} y \texttt{monto}, lee \textbf{solo esas dos} de las 11 columnas.
\end{block}
\centering
\begin{tikzpicture}[font=\scriptsize]
  \foreach \cname [count=\i] in {pedido\_id,fecha,cliente\_id,distrito,categoria,producto,cantidad,precio,monto,metodo\_pago,estado} {
    \pgfmathsetmacro{\x}{(\i-1)*0.55}
    \ifnum\i=5 \def\col{utpRojo!70}\else\ifnum\i=9 \def\col{utpRojo!70}\else\def\col{utpGrisClaro}\fi\fi
    \fill[\col] (\x,0) rectangle (\x+0.42,2.4);
    \draw[utpGrisMedio] (\x,0) rectangle (\x+0.42,2.4);
    \node[rotate=90,anchor=east] at (\x+0.21,-0.1) {\cname};
  }
\end{tikzpicture}
\column{.49\textwidth}
\begin{block}{2. Reparte el trabajo entre muchas máquinas}
\small Cada máquina suma su parte al mismo tiempo; al final se juntan los resultados.
\end{block}
\centering
\begin{tikzpicture}[font=\scriptsize,
 caja/.style={draw=utpRojo,thick,rounded corners=3pt,align=center}]
  \node[caja,fill=utpRojo!25,minimum width=5.6cm] (t) at (0,2.2) {\footnotesize 1.2 millones de filas};
  \foreach \i in {1,...,4} {
    \node[caja,fill=white,minimum width=1.2cm,minimum height=0.9cm] (m\i) at (-2.8+\i*1.4,0.9) {máquina \i\\300 mil};
    \draw[-latex] (t.south -| m\i.north) -- (m\i.north);
  }
  \node[caja,fill=utpRojo!10,minimum width=3.2cm] (r) at (0,-0.4) {\footnotesize suma de las 4 partes};
  \foreach \i in {1,...,4} \draw[-latex] (m\i.south) -- (r.north);
\end{tikzpicture}\par\vspace{2mm}
{\footnotesize Google usa cientos o miles de máquinas. Tú solo escribes la consulta.}
\end{columns}
\end{frame}

% 14
\separador{Parte 2: BigQuery}{Guardar tablas enormes y hacerles preguntas con SQL}

% 15
\begin{frame}{¿Qué es BigQuery?}
\begin{block}{Definición}
\small Servicio de Google Cloud para guardar tablas muy grandes y responder preguntas sobre ellas con SQL. \textbf{No creas servidores:} cargas los datos y escribes la consulta.
\end{block}
\renewcommand{\arraystretch}{1.15}
\centering\small
\begin{tabular}{@{}p{2.6cm}>{\raggedright\arraybackslash}p{4.9cm}>{\raggedright\arraybackslash}p{5.0cm}@{}}
\toprule
 & \textbf{Cloud SQL} & \textbf{BigQuery}\\
\midrule
\textbf{Para qué} & Registrar lo que pasa en la app & Analizar el historial\\
\textbf{Qué creas} & Una instancia: eliges motor y máquina & Un conjunto de datos y sus tablas; ningún servidor\\
\textbf{Cuántas filas} & Miles o millones & Millones o miles de millones\\
\textbf{Consulta típica} & Trae el pedido P-1001 & Suma las ventas de 3 años por categoría\\
\textbf{Cómo cobra} & Por las horas que la instancia está encendida & Por los datos guardados y los que lee cada consulta\\
\bottomrule
\end{tabular}
\par\vspace{2mm}
\begin{alertblock}{Las dos hablan SQL}
\small Lo que aprendas a consultar en una sirve en la otra.
\end{alertblock}
\end{frame}

% 16
\begin{frame}[fragile]{Cómo se organiza BigQuery}
\begin{center}
\begin{tikzpicture}[font=\small,
 nivel/.style={draw=utpRojo,thick,rounded corners=4pt,align=center,minimum width=4.4cm,minimum height=0.75cm}]
  \node[font=\bfseries] at (0,2.9) {Cloud SQL};
  \node[nivel,fill=utpRojo!30] (a1) at (0,2.2) {Instancia \texttt{mysql-server}};
  \node[nivel,fill=utpRojo!14] (a2) at (0,1.15) {Base de datos \texttt{ventas}};
  \node[nivel,fill=white] (a3) at (0,0.1) {Tabla \texttt{pedidos}};
  \draw[-latex,thick] (a1) -- (a2); \draw[-latex,thick] (a2) -- (a3);
  \node[font=\bfseries] at (7.6,2.9) {BigQuery};
  \node[nivel,fill=utpRojo!30] (b1) at (7.6,2.2) {Proyecto de Google Cloud};
  \node[nivel,fill=utpRojo!14] (b2) at (7.6,1.15) {Conjunto de datos \texttt{billing\_dataset}};
  \node[nivel,fill=white] (b3) at (7.6,0.1) {Tabla \texttt{sampleinfotable}};
  \draw[-latex,thick] (b1) -- (b2); \draw[-latex,thick] (b2) -- (b3);
  \foreach \y/\t in {1.15/grupo de tablas,0.1/donde están las filas} {
    \node[font=\scriptsize,color=utpGris] at (3.8,\y) {\t};
  }
\end{tikzpicture}
\end{center}
\begin{columns}[T,onlytextwidth]
\column{.48\textwidth}
\begin{block}{En la consulta: \texttt{conjunto.tabla}}
\begin{lstlisting}
SELECT *
FROM `billing_dataset.sampleinfotable`
LIMIT 10
\end{lstlisting}
\end{block}
\column{.48\textwidth}
\begin{exampleblock}{En la consola de BigQuery}
\small En el panel \textbf{Explorador}, los tres puntos (\textbf{Ver acciones}) junto al proyecto llevan a \textbf{Crear conjunto de datos}; junto al conjunto de datos, a \textbf{Crear tabla}.
\end{exampleblock}
\end{columns}
\end{frame}

% 17
\begin{frame}[fragile]{SQL para analizar (1): elegir y filtrar}
\begin{columns}[c,onlytextwidth]
\column{.28\textwidth}
\centering
{\small\bfseries Tabla \texttt{pedidos}}\par\vspace{1mm}
\scriptsize\setlength\tabcolsep{3pt}
\begin{tabular}{@{}llr@{}}
\toprule
\textbf{pedido} & \textbf{categoria} & \textbf{monto}\\
\midrule
\filaA P-1 & Moda & 98\\
\filaA P-2 & Electrónica & 295\\
P-3 & Supermercado & 12\\
P-4 & Supermercado & 0\\
P-5 & Supermercado & 56\\
P-6 & Electrónica & 59\\
\bottomrule
\end{tabular}
\column{.04\textwidth}
\column{.34\textwidth}
\begin{lstlisting}
SELECT pedido, monto
FROM pedidos
WHERE monto > 60
ORDER BY monto DESC
\end{lstlisting}
\centering
\begin{tikzpicture}\draw[-latex,very thick,utpRojo] (0,0) -- (1.2,0);\end{tikzpicture}\par\vspace{1mm}
\footnotesize
\begin{tabular}{@{}lr@{}}
\toprule
\textbf{pedido} & \textbf{monto}\\
\midrule
P-2 & 295\\
P-1 & 98\\
\bottomrule
\end{tabular}
\column{.28\textwidth}
\renewcommand{\arraystretch}{1.25}
\footnotesize
\begin{tabular}{@{}>{\color{utpRojo}\bfseries\ttfamily}l>{\raggedright\arraybackslash}p{2.0cm}@{}}
SELECT & qué columnas\\
FROM & de qué tabla\\
WHERE & qué filas\\
ORDER BY & en qué orden\\
LIMIT & cuántas filas\\
\end{tabular}
\end{columns}
\par\vspace{4mm}
\begin{alertblock}{Se lee así}
\small ``Muéstrame el pedido y el monto de la tabla \texttt{pedidos}, solo los que pasan de 60, del más caro al más barato''.
\end{alertblock}
\end{frame}

% 18
\begin{frame}[fragile]{SQL para analizar (2): agrupar, contar y sumar}
\begin{columns}[c,onlytextwidth]
\column{.28\textwidth}
\centering
{\small\bfseries Tabla \texttt{pedidos}}\par\vspace{1mm}
\scriptsize\setlength\tabcolsep{3pt}
\begin{tabular}{@{}llr@{}}
\toprule
\textbf{pedido} & \textbf{categoria} & \textbf{monto}\\
\midrule
\filaB P-1 & Moda & 98\\
\filaA P-2 & Electrónica & 295\\
\filaC P-3 & Supermercado & 12\\
\filaC P-4 & Supermercado & 0\\
\filaC P-5 & Supermercado & 56\\
\filaA P-6 & Electrónica & 59\\
\bottomrule
\end{tabular}
\column{.04\textwidth}
\column{.34\textwidth}
\begin{lstlisting}
SELECT categoria,
  COUNT(*) AS pedidos,
  SUM(monto) AS total
FROM pedidos
GROUP BY categoria
ORDER BY total DESC
\end{lstlisting}
\column{.30\textwidth}
\centering
{\small\bfseries Resultado}\par\vspace{1mm}
\scriptsize\setlength\tabcolsep{3pt}
\begin{tabular}{@{}lrr@{}}
\toprule
\textbf{categoria} & \textbf{pedidos} & \textbf{total}\\
\midrule
\filaA Electrónica & 2 & 354\\
\filaB Moda & 1 & 98\\
\filaC Supermercado & 3 & 68\\
\bottomrule
\end{tabular}
\end{columns}
\par\vspace{4mm}
\begin{block}{Qué hace cada parte}
\small \texttt{GROUP BY} junta las filas con la misma categoría (mismo color). \texttt{COUNT(*)} cuenta las filas de cada grupo y \texttt{SUM(monto)} suma sus montos.
\end{block}
\begin{alertblock}{Mira el resultado}
\small Supermercado tiene \textbf{más pedidos}, pero Electrónica \textbf{vende más dinero}. Contar y sumar responden preguntas distintas.
\end{alertblock}
\end{frame}

% 19
\begin{frame}[c]{Práctica 1}
\begin{block}{Realiza el laboratorio}
\large \textit{Examina los datos de facturación con BigQuery}\par\vspace{1mm}
\small\texttt{skills.google/catalog\_lab/1784}
\end{block}
\begin{exampleblock}{Al terminar}
\small Sabes qué servicio de Google Cloud generó el mayor costo en los datos de ejemplo.
\end{exampleblock}
\end{frame}

% 20
\begin{frame}[fragile]{Práctica 2: BigQuery sin BigQuery}
\begin{columns}[T,onlytextwidth]
\column{.44\textwidth}
\begin{block}{Qué harás}
\small Consultas como las del lab con \textbf{Python} y \textbf{DuckDB}: un motor SQL que corre dentro de Python y, como BigQuery, lee por columnas. No necesitas cuenta de Google Cloud.
\end{block}
\begin{block}{De dónde salen los datos}
\small No se descargan: la primera celda del notebook ejecuta \texttt{generar\_ventas.py}, que crea \texttt{ventas.csv} con \textbf{1.2 millones} de pedidos de Pedidos Express (2023--2025). Todos obtienen los mismos datos.
\end{block}
\column{.53\textwidth}
\begin{lstlisting}
import duckdb

duckdb.sql("""
  SELECT categoria,
    ROUND(SUM(monto), 2) AS total_soles
  FROM ventas
  GROUP BY categoria
  ORDER BY total_soles DESC
""")
\end{lstlisting}
\par\vspace{2mm}
\begin{exampleblock}{Cómo empezar}
\scriptsize 1. En Colab: \textbf{Archivo} $\rightarrow$ \textbf{Abrir cuaderno} $\rightarrow$ \textbf{GitHub}.\par
2. Pega el repositorio del curso:\par
\hspace*{3mm}\texttt{github.com/rgap/}\par
\hspace*{3mm}\texttt{rgap-Integracion-de-Servicios-Cloud-Gcp-Agosto2026}\par
3. Elige \texttt{Sesion24/practica/practica\_bigquery\_python.ipynb}, ejecuta las celdas en orden y resuelve los 5 retos del final.
\end{exampleblock}
\end{columns}
\end{frame}

% 21
\begin{frame}{Cierre: preguntas para comprobar}
\begin{block}{Responde con tus palabras}
\small
1. En el caso de estudio, ¿por qué basta una instancia con cuatro bases de datos?\\\medskip
2. Nombra las tres V de Big Data con un ejemplo de Pedidos Express.\\\medskip
3. ¿Por qué BigQuery lee menos datos si la consulta usa solo dos columnas?\\\medskip
4. En una consulta que suma los montos por categoría, ¿qué hace \texttt{GROUP BY}?\\\medskip
5. En la práctica, ¿la categoría con más pedidos fue la que más dinero vendió? ¿Por qué?
\end{block}
\begin{alertblock}{Idea final}
\small Cloud SQL registra el día a día; BigQuery responde preguntas sobre todo el historial. Las dos se consultan con SQL.
\end{alertblock}
\end{frame}

% 22
\begin{frame}{Referencias}
\scriptsize
\begin{itemize}\setlength\itemsep{4pt}
  \item Laboratorio: \textit{Examina los datos de facturación con BigQuery} (\textit{Examining Billing Data with BigQuery}), curso \textit{Essential Google Cloud Infrastructure: Core Services}: \texttt{skills.google/catalog\_lab/1784}.
  \item Práctica de la sesión: \texttt{github.com/rgap/rgap-Integracion-de-Servicios-Cloud-Gcp-Agosto2026}, carpeta \texttt{Sesion24/practica/} (\texttt{practica\_bigquery\_python.ipynb} y \texttt{generar\_ventas.py}).
  \item DuckDB: \texttt{duckdb.org}.
  \item Material previo del curso: Sesión 23, bases de datos en Google Cloud y Cloud SQL.
  \item Google Cloud: BigQuery, Cloud SQL y Cloud Storage.
  \item Las respuestas de las preguntas están en \texttt{../respuestas.md}.
\end{itemize}
\end{frame}

\end{document}
